% Figure — Exercice « Où est la contre-réaction ? » : montage (a) et schéma équivalent aux petites variations en quadripôles A et β_F (b) (EANA3, chapitre 2).
% (a) D'après le support [support, p. 11, figure 1], redessiné.
% (b) D'après les notes externes [notes-2023-c08, p. 2-3] et le corrigé externe [corr-2023-recueil, p. 3], redessiné.
% Rendu SVG :
%   latex -output-directory=<dossier temporaire> 02-ou-est-la-contre-reaction.tex
%   dvisvgm --no-fonts --exact-bbox -o 02-ou-est-la-contre-reaction.svg <dossier temporaire>/02-ou-est-la-contre-reaction.dvi
\documentclass[border=4pt]{standalone}
\usepackage[T1]{fontenc}
\usepackage[utf8]{inputenc}
\usepackage{amsmath}
\usepackage{sansmath}
\renewcommand{\familydefault}{\sfdefault}
\usepackage[european, straightvoltages, americancurrents]{circuitikz}
\begin{document}
\sansmath
\begin{circuitikz}[font=\small, >=latex, fleche/.style={->, thick}]
  % ---------- (a) ----------
  \node at (-0.4,5.4) {(a)};
  \draw (4.0,1.6) node[npn, anchor=B] (Q) {};
  \node[right] at (Q.text) {$M_1$};
  \draw (Q.E) to[R, l=$R_x$] (Q.E |- 0,0);
  \draw (Q.C) -- ++(0,0.4) coordinate (c) to[R, l=$R_C$] ++(0,2.0) coordinate (vccC);
  \draw (2.6,1.6) to[R, l=$R_b$] (2.6,4.4 |- vccC) -- (vccC);
  \node[above] at (3.8,4.4 |- vccC) {$V_{CC}$};
  \draw (3.8,4.4 |- vccC) -- ++(0,0.2);
  \draw (0,1.6) node[ocirc]{} -- (1.0,1.6) to[C, l=$C_1$] (2.6,1.6) -- (Q.B);
  \draw (1.0,1.6) to[R, l_=1 k$\Omega$] (1.0,0);
  \draw (0,0) node[ocirc]{} -- (6.2,0) node[ocirc]{};
  \draw (c) -- (6.2,0 |- c) node[ocirc]{};
  \draw[fleche] (-0.3,0.15) -- node[left]{$V_i$} (-0.3,1.45);
  \draw[fleche] (6.5,0.15) -- node[right]{$V_o$} (6.5,0 |- c);
  % ---------- (b) ----------
  \begin{scope}[xshift=8.6cm]
    \node at (-0.4,5.4) {(b)};
    \draw[thick, fill=black!3] (0.9,2.0) rectangle (6.2,4.9);  \node[font=\bfseries] at (5.9,4.65) {$A$};
    \draw[thick, fill=black!3] (0.9,-0.6) rectangle (6.2,1.3); \node at (5.5,-0.3) {$\beta_F$};
    \draw (0,4.2) node[ocirc]{} -- (1.8,4.2) to[R, l=$r_{be}$] (1.8,2.5) -- (3.4,2.5);
    \draw[fleche] (0.4,4.45) -- node[above]{$i_b$} (1.3,4.45);
    \draw[fleche] (1.25,2.65) -- node[left]{$v_{be}$} (1.25,4.05);
    \draw (3.4,4.2) to[american controlled current source, l=$g_m v_{be}$] (3.4,2.5);
    \draw (5.3,4.2) to[R, l=$r_{ce}$] (5.3,2.5);
    \draw (3.4,4.2) -- (7.2,4.2) node[ocirc]{};  \draw (3.4,2.5) -- (5.3,2.5);
    \draw (2.6,2.5) node[circ]{} -- (2.6,1.0) -- (3.3,1.0) to[R, l=$R_x$] (3.3,0) -- (3.3,-1.2);
    \draw[fleche] (2.85,0.85) -- node[left]{$i_o$} (2.85,0.15);
    \draw (0,-1.2) node[ocirc]{} -- (7.2,-1.2) node[ocirc]{};
    \draw (0,-1.2) node[ground]{};
    \draw[fleche] (-0.3,-1.0) -- node[left]{$v_i$} (-0.3,4.0);
    \draw[fleche] (7.0,4.45) -- node[above]{$i_o$} (6.4,4.45);
    \draw[fleche] (7.5,-1.0) -- node[right]{$v_o$} (7.5,4.0);
  \end{scope}
\end{circuitikz}
\end{document}
